\section{Integral ray quotients through a split irregular CM point}
\label{sec:cubic-irregular-ray}

The ray moduli in this section are squarefree ideals prime to a fixed
set containing \(6p(269)\).  This is the set of auxiliary moduli
needed for the rank-two Euler-system argument.  No assertion about
arbitrary repeated auxiliary prime powers is needed.

\begin{lemma}[AC-K29-irregular-old: saturation at the displayed ray level]
\label{AC-K29-irregular-old}
Let \(p\geq5\) split in \(K=\mathbb Q(\omega)\), let
\(\mathfrak n\) be a permitted squarefree auxiliary ideal,
and let \(\nu\) be a finite \(p\)-power ray character of conductor
dividing \(\mathfrak n\).  Put \(\psi=\chi^a\nu\), for
\(a=1,2\), and let its primitive theta series have tame level
\(M\).  Use the displayed tame level
\(3\operatorname N((269)\mathfrak n)\), with its theta
stabilizations or depletions.  Complete the full ordinary Hecke algebra in
characteristic zero at the resulting weight-one theta eigenform,
with a chosen ordinary refinement at \(p\).  After a strict
coefficient extension, take either of its CM branch maps to
\(A=E[[X]]\).

In the corresponding actual cohomology quotient, Ohta's ordinary
line maps isomorphically to its rank-one character line in the
torsionfree induced quotient.  This includes the case of repeated
ordinary roots at \(p\) and repeated tame stabilization roots.
\end{lemma}

\begin{proof}
We first justify the precise raised-level algebra used in the proof.
The primary conductor input is Dimitrov,
\emph{On the local structure of ordinary Hecke algebras at classical
weight one points}, Proposition 6.5, printed pp.~11--12,
\url{https://gitlabpages.univ-lille.fr/dimitrov/articles/durham.pdf}.
It states, without a regularity hypothesis at \(p\), that all
classical weight-one specializations of a primitive Hida family
have its tame level and local type.  Its complete proof in \S6
uses finite inertia and excludes a special family by the
Frobenius-eigenvalue ratio \(q\).  Thus every primitive family
through the raised-level point has the primitive tame level \(M\).
In particular no family new at an extra tame prime passes through
that point.

Let \(T_0\) be the full completed ordinary algebra at the primitive
theta series, and put \(D_q=q\langle q\rangle\) in the universal
covariant convention.  If the exponent of an extra rational prime
\(q\neq p\) is at most two, the relevant oldform algebras are
the local factors, at the selected \(U_q\)-value, of
\[
\begin{array}{c|c}
\text{primitive/displayed exponents}&\text{monic polynomial in }Y\cr
0/1&Y^2-T_qY+D_q\cr
0/2&Y(Y^2-T_qY+D_q)\cr
1/2&Y(Y-U_q).
\end{array}                                                    \tag{CIR.1}
\]
There is no extra polynomial when the two exponents agree.
The last row occurs only when the primitive local coefficient
\(U_q\) is nonzero.  For the squarefree ideal moduli at issue,
a nontrivial character at one of two split primes has precisely
this ramified-principal-series local type.  A character ramified
at both split primes, or at an inert prime, already has exponent
two.  Characters of smaller conductor are covered by the first
two rows.  The fixed cubic conductor \((269)\) cannot cancel
against an auxiliary \(p\)-power character, since \(p\neq3,269\).

Here is an algebraic verification that (CIR.1) describes the actual
completed Hecke algebra, including repeated roots.  On a classical
old space the vectors \(f(z),f(qz),\ldots\) are independent, as is
seen successively from their first nonzero Fourier coefficients.
The \(U_q\)-operator is cyclic on that space, with the displayed
characteristic and minimal polynomial; this follows either by
applying \(U_q\) to the Fourier series, or from the companion
matrix in the proof of LLZ15, Theorem 4.2.8.  Different primes
give tensor products of these cyclic spaces.

The following density arguments are made on finite affinoid
ordinary Hecke algebras before taking formal completions.  Shrink
the weight neighborhood and the finitely many Hecke neighborhoods
around the points in question so that primitive components not
passing through those points are excluded.  All relations being
checked have finite presentations on these affinoids.  Positive
classical weights are dense there by ordinary control.  Complete
only after proving the finite-algebra identities.  Thus density
is not being applied to an arbitrary, possibly nonconvergent
formal power series in \(E[[X]]\).

The lower-level Hecke operators belong to the completed
higher-level algebra.  Operators away from the raised level and
\(U_p\) already do.  An omitted good-prime \(T_q\) is recovered
as the trace of Frobenius on the rank-two Galois representation;
at a prime in \(M\), the primitive \(U_q\) is Frobenius on the
inertia-invariant line, or is zero when that invariant space is
zero.  These assertions are valid on the completed ring, not just
its generic points.  Indeed the proof of Bellaiche--Dimitrov,
arXiv:1301.0712v2, Proposition 7.1, printed pp.~19--20, gives
finite inertia near an Artin point.  In the nonreduced case the
same argument uses the invertibility of
\(\operatorname{Ad}(\mathrm{Frob}_q)-q\) at the point to kill
the monodromy operator: its residual eigenvalues are roots of
unity, and none equals \(q\).  A finite inertia representation
over a characteristic-zero complete local ring has the same
invariant rank as its residue representation, by averaging.
This also gives the indicated Frobenius operator on its invariant
line.  The monodromy input is exactly the family monodromy theorem
used in that printed proof; no smoothness or regularity-at-\(p\)
conclusion from that paper is invoked.

One can also first work in the subalgebra generated by the
good-prime operators and \(U_p\).  It is reduced, since these
operators are semisimple on every dense classical old space;
the Galois pseudorepresentation takes values there.  The
finite-inertia conclusion on that reduced algebra then extends
to the full old algebra, including its possible nilpotents.
The new tame diamond characters have finite order, so their
values are fixed in the characteristic-zero completion and
add no infinitesimal parameter.

Relations among the resulting lower-level operators hold in the
higher-level algebra.  To check this, specialize to the dense
classical weights at least two and use the primitive oldform
decomposition.  The lower-level operators act as scalars on each
primitive block, so this argument also kills the entire old
Jordan block, not merely its eigenvalues.  Ordinary control and
the torsionfreeness over the weight algebra then give the relation
before specialization.  We obtain a surjection from the local
factor of the tensor product of the algebras (CIR.1) over \(T_0\)
onto the actual completed algebra \(T\).
There is no kernel.  Hensel-factor each polynomial into the
selected monic factor at the chosen residual \(U_q\)-value and
its coprime complement.  On a sufficiently small affinoid
neighborhood the same factorization and spectral summands are
defined.  Reduce a proposed relation successively modulo these
selected monic factors.  On the corresponding summand of every
classical primitive old block, cyclicity of the selected
companion matrices makes every remaining
coefficient zero.  Density and torsionfreeness on \(T_0\) make
those coefficients zero in \(T_0\).  This proves the claimed
algebra description, including the completed local factors.
The primitive oldform decomposition of ordinary families used
here is the one cited and applied in Dimitrov, Lemma 6.1 and its
complete proof, printed pp.~9--10; the companion degeneracy
calculation is LLZ15, Definition 4.2.7 and Theorem 4.2.8.

Each monic extension \(C=B[Y]/(P)\) is a finite free
Gorenstein extension of \(B\).  Explicitly, the functional
extracting the coefficient of \(Y^{\deg P-1}\) gives a perfect
\(B\)-pairing \((u,v)\mapsto\lambda(uv)\) on \(C\).
Its matrix is triangular about the anti-diagonal with unit
anti-diagonal entries.  Hence
\(\operatorname{Hom}_B(C,B)\simeq C\).
The same assertion holds for each completed local factor and
for successive tensor products.  Consequently, for the common
weight algebra \(\Lambda\),
\[
 \omega_T=\operatorname{Hom}_{\Lambda}(T,\Lambda)
       \simeq T\otimes_{T_0}\omega_{T_0},\qquad
 \omega_T\otimes_T A\simeq\omega_{T_0}\otimes_{T_0}A.  \tag{CIR.2}
\]

If the ordinary roots at \(p\) coincide, the right side is
\(A\oplus k^2\) by AC-K29-irregular-algebra.  The exact same
cohomological argument as in AC-K29-irregular-saturation applies
at level \(T\): Ohta gives an injective line \(A\) with this
cokernel; global character separation gives two copies of a
rank-one \(A\)-module; and the torsion-length equation
\(a+2b=2\) forces \(a=0\).  Thus the free-coordinate index of
the ordinary line is zero.

If the ordinary roots are distinct, the conclusion is simpler
and needs no description of the minimal algebra.  On the branch
let the ordinary line and unramified quotient characters of
Frobenius be \(\alpha(X),\beta(X)\), with
\(\alpha(0)\neq\beta(0)\).  The operator
\((\mathrm{Frob}_p-\beta)/(\alpha-\beta)\) maps the cohomology
quotient to the ordinary line and restricts there to the identity.
It splits that line.  Its image in the torsionfree quotient
therefore has index zero as well.  This proves the lemma.
\end{proof}

\begin{proposition}[AC-K29-irregular-ray: integral repeated-root ray quotients]
\label{AC-K29-irregular-ray}
Let \(p\geq5\), \(p\neq269\), split in
\(K=\mathbb Q(\omega)\), and let \(\chi\) be the cubic character
of conductor \((269)\).  No distinction between the two values
\(\chi(\mathfrak p),\chi(\overline{\mathfrak p})\) is required.
For \(a=1,2\), either choice of CM ordinary branch, and every
permitted squarefree auxiliary modulus \(\mathfrak n\), there
are actual integral cohomological quotient maps to
\[
 \operatorname{Ind}_{G_{K(\mathfrak n)}}^{G_{\mathbb Q}}
                   \mathcal O(\chi^{-a})               \tag{CIR.3}
\]
at weight one, simultaneously compatible with the LLZ ray norm
maps.  Their denominators are independent of \(\mathfrak n\);
the CM-side maps themselves are integral.  At the base modulus
their two ordinary lines are the two different global
\(\chi^{-a}\) and \(\chi^{a}\) lines, even when local Frobenius
at \(p\) is scalar.
\end{proposition}

\begin{proof}
\proofstep{1. Explicit integral coefficient rings and actual modules.}
Write \(\Gamma_{\mathfrak n}=\operatorname{Gal}(K(\mathfrak n)/K)\)
for the maximal abelian \(p\)-extension of conductor dividing
\(\mathfrak n\).  Put
\[
 R_{\mathfrak n}=\mathcal O[\Gamma_{\mathfrak n}][[X]].  \tag{CIR.4}
\]
The field \(K\) has class number one.  Since \(p\nmid6\),
the pro-\(p\) local units at \(\mathfrak p\) generate its
\(\mathfrak p\)-ramified \(\mathbb Z_p\)-extension.
It is disjoint from the ray extensions in (CIR.4), which are
unramified at \(p\).  Choose an inertia element \(s\) mapping
to a generator of that \(\mathbb Z_p\)-extension.  Use the
cohomological character convention so that the inducing
character \(\Theta_{\mathfrak n}\) satisfies
\[
 \Theta_{\mathfrak n}(s)=1+X,\qquad
 \Theta_{\mathfrak n}^{c}(s)=1.                        \tag{CIR.5}
\]
Changing to the inverse generator only changes the coordinate
by a unit times \(X\).

Let \(\Psi_{\mathfrak n}\) denote the automorphic universal
Hecke character; its cohomological contragredient is
\(\Theta_{\mathfrak n}\).  Thus the weight coordinate in
(CIR.5) is chosen on the latter character, and the automorphic
inertia value is \((1+X)^{-1}\).
The usual integral theta Fourier series, with coefficients
\(\sum_{\operatorname N\mathfrak a=m}
\Psi_{\mathfrak n}(\mathfrak a)\), where
\(\Psi_{\mathfrak n}(\mathfrak a)
=\chi^a(\mathfrak a)[\mathfrak a]
\Psi_{\mathfrak p}(\mathfrak a)\),
and the ideals dividing the displayed modulus omitted, defines
an ordinary Hecke map to \(R_{\mathfrak n}\).
This assertion does not use a cohomological freeness theorem:
evaluate at finite characters of \(\Gamma_{\mathfrak n}\)
and the dense positive classical weights of the CM family.
Each evaluation is the usual theta eigenform; hence every
Hecke relation holds.  The group ring is torsionfree and its
characteristic-zero character evaluations separate its elements.
This proves the integral map, with the ordinary unit root
given by the ideal above the other prime of \(K\).
The construction of the CM family itself is precisely the
theta series in BD21, \S3.1, before Proposition 3.2.
The two branches exchange \(\mathfrak p\) and
\(\overline{\mathfrak p}\).

Let \(\mathcal H_{\mathfrak n}\) denote the full ordinary
cohomology, localized at the common non-Eisenstein residual
theta ideal, and set
\(J_{\mathfrak n}=R_{\mathfrak n}\otimes_{\mathbb T_{\mathfrak n}}
\mathcal H_{\mathfrak n}\).
Ohta's ordinary sequence base changes to
\[
 R_{\mathfrak n}\xrightarrow{\iota_{\mathfrak n}}
 J_{\mathfrak n}\longrightarrow Q_{\mathfrak n}
 \longrightarrow0,                                   \tag{CIR.6}
\]
where inertia at \(p\) is trivial on \(Q_{\mathfrak n}\), and
\(s\) acts by \(1+X\) on the source.  The first map is
injective.  It is injective after passing to the generic field
of the weight algebra and every finite ray character, by
classical control and the ordinary line of the theta series;
its source is free over \(\mathcal O[[X]]\), so it has no
element killed by this passage.  Notice that no freeness of
\(J_{\mathfrak n}\) is asserted.

\proofstep{2. A uniform coordinate extraction and actual division.}
There is a unique \(R_{\mathfrak n}\)-linear map
\(\varphi_{\mathfrak n}:J_{\mathfrak n}\to R_{\mathfrak n}\)
such that
\[
 \iota_{\mathfrak n}\varphi_{\mathfrak n}(h)=(s-1)h.
\]
It exists by (CIR.6).  Define
\[
 F_{\mathfrak n}(h)=
       (\varphi_{\mathfrak n}(h),\varphi_{\mathfrak n}(ch)).
                                                               \tag{CIR.7}
\]
At a generic CM fiber this is \(X\) times the two induced
coordinates normalized by the ordinary line.  Therefore it is
Galois equivariant as a map to the induced representation over
\(R_{\mathfrak n}\): the identities hold at the generic fibers
and the target is torsionfree over \(\mathcal O[[X]]\).
The same argument gives exactly
\[
 F_{\mathfrak n}\iota_{\mathfrak n}(1)=(X,0),\qquad
 F_{\mathfrak n}c\iota_{\mathfrak n}(1)=(0,X).           \tag{CIR.8}
\]

We now prove the essential divisibility, not merely integrality.
For every finite ray character \(\nu\), pass to its coefficient
field and complete at \(X=0\).  Its theta series has primitive
character \(\chi^a\nu\); this remains nontrivial on conjugation,
since its cubic residual character is nontrivial.  The preceding
old-level lemma says that its actual ordinary line has index
zero in the torsionfree induced quotient.  Thus the coordinate
map (CIR.7) has image in \(XA^2\), including on torsion, which
maps to zero in the torsionfree target.  It follows that
\(F_{\mathfrak n}\bmod X\) vanishes after every ray character
and after inverting \(p\).  But
\[
 R_{\mathfrak n}/XR_{\mathfrak n}
                 =\mathcal O[\Gamma_{\mathfrak n}]      \tag{CIR.9}
\]
is \(\mathcal O\)-torsionfree, and its characteristic-zero
character evaluations are faithful.  Hence the same reduction
is zero integrally.  We have proved
\[
 F_{\mathfrak n}(J_{\mathfrak n})\subset XR_{\mathfrak n}^2,
 \qquad \pi_{\mathfrak n}=X^{-1}F_{\mathfrak n}.
                                                               \tag{CIR.10}
\]
This is a well-defined integral Galois-equivariant quotient map.
Its surjectivity follows from (CIR.8).  Its kernel is saturated
over \(\mathcal O[[X]]\), since its target is free over that
ring.  Formula (CIR.10) has no ray-character idempotent
denominators.

The ordinary Hecke map itself may also be seen to be onto.
Choose \(h_0\in G_K\) with cubic value \(\omega\) and trivial
image in the maximal abelian pro-\(p\) quotient.  The cubic and
pro-\(p\) quotients are disjoint.  If \(t_g\) is the image of
the universal Galois trace, then
\[
 \Theta_{\mathfrak n}(g)
       ={t_{h_0g}-\omega^{a}t_g\over\omega^{-a}-\omega^{a}}
                                                        \tag{CIR.11}
\]
in the convention where the first inducing character has cubic
part \(\chi^{-a}\).  The denominator is a fixed unit.  The weight
algebra already supplies the \(\mathfrak p\)-character values;
dividing those units recovers every ray group element.  Thus the
construction can be regarded as a literal cohomological quotient,
not only a scalar extension followed by a map.

\proofstep{3. Surjective actual norm maps without a freeness assumption.}
The norm operators are precisely LLZ15, Definition 3.3.1,
with the universal theta coefficients.  Their Hecke equivariance
is Proposition 3.3.2 and its complete elementary calculation.
We need their surjectivity, not the injectivity that was deduced
from \(p\)-distinguished freeness in the source.

For stabilization, the proof of LLZ Theorem 4.2.8 first uses
Ihara's surjection and Lemma 4.2.4 to obtain a surjection onto
two copies of the lower cohomology.  Taking the combination
\((u,v)\mapsto u-(T_q-Y)v/q\) is still surjective.  This is
exactly its map \(\beta\).  None of these surjectivity steps
uses the freeness hypothesis employed earlier in that proof
to prove injectivity.  For depletion, the entire surjectivity
argument of Theorem 4.2.9 uses the exact degeneracy sequence
of Lemma 4.2.2 and the generation by \(U_q\); again it does
not use that freeness hypothesis.  The inert map is the
composition giving Corollary 4.2.10.  All auxiliary primes
are different from the coefficient prime \(p\).

These surjections pass to full ordinary cohomology.  Explicitly,
form the corresponding maps of finite modules over the ordinary
weight algebra, reduce modulo a finite ordinary control ideal,
and apply those same finite-level surjections.  Ordinary control
holds before a \(p\)-distinguished localization.  Nakayama then
gives surjectivity of the big maps.  Only finite generation is
used here; a claim that their full sources are free or that the
big maps are isomorphisms would be unwarranted.  Finally tensor
by the theta maps and the quotient
\(R_{\mathfrak n'}\to R_{\mathfrak n}\).  Right exactness gives
the needed surjections
\[
 \beta_{\mathfrak n',\mathfrak n}:
 R_{\mathfrak n}\otimes_{R_{\mathfrak n'}}J_{\mathfrak n'}
                      \twoheadrightarrow J_{\mathfrak n}.       \tag{CIR.12}
\]
The maps preserve Ohta's ordinary submodule.  Besides the
functorial connected-part construction, this follows directly
before coefficient base change: the composite from the source
ordinary line into the target unramified quotient is killed
by the nonzero ordinary inertia parameter, whereas that target
quotient is free over the weight algebra by Ohta's duality.
The composite is therefore zero.  This assertion then survives
base change.  Write the resulting rank-one scalar as
\[
 \beta_{\mathfrak n',\mathfrak n}\iota_{\mathfrak n'}(1)
                 =r_{\mathfrak n',\mathfrak n}
                         \iota_{\mathfrak n}(1).        \tag{CIR.13}
\]
No unit property is assumed in this formula.

The norm maps commute with \(s\) and \(c\).  Applying
\(s-1\), and using the injectivity in (CIR.6), proves
\(\varphi_{\mathfrak n}\beta
 =r_{\mathfrak n',\mathfrak n}\varphi_{\mathfrak n'}\).
Applying the same identity to \(ch\), and then dividing the
actual identity by \(X\), proves
\[
 \pi_{\mathfrak n}\beta_{\mathfrak n',\mathfrak n}
       =r_{\mathfrak n',\mathfrak n}\pi_{\mathfrak n'}.
                                                               \tag{CIR.14}
\]
Both \(\beta\) and \(\pi_{\mathfrak n}\) are surjective.
Since \(\pi_{\mathfrak n'}\) is also onto, (CIR.14) implies
\(r_{\mathfrak n',\mathfrak n}R_{\mathfrak n}=R_{\mathfrak n}\).
The scalar is therefore a unit.  This proves its needed
normalization directly; no unverified Fricke or differential
normalization is used to declare it a unit.

Choose the ordinary-line generators compatibly along an
increasing cofinal sequence of squarefree products of allowed
prime ideals.  A unit lifts under the surjection of the finite
local group rings in (CIR.4); the commuting norm diagrams then
give compatible choices for all intermediate moduli, exactly
as in the unit-patching proof of LLZ, Corollary 5.2.6.
With these choices every scalar in (CIR.14) is one.  Specializing
\(X=0\) proves (CIR.3) and its simultaneous integral norm
compatibility.

For \(p=67\) the fixed conductor ray group does have a
\(67\)-part, since \((269^2-1)/6=12060\).  The explicit ring
(CIR.4) varies only the auxiliary ray group; the fixed conductor
character is already \(\chi^a\).  We set the fixed \(p\)-part
to the trivial character by a quotient map, not by averaging
over that group.  Thus neither that fixed factor nor the
growing auxiliary groups introduce an unbounded denominator.
Finally, the two branches at the base modulus have ordinary
lines distinguished by the global cubic separator \(h_0\).
Their local Frobenius being scalar does not identify these
two one-dimensional subspaces.  This completes the proof.
\end{proof}

\begin{theorem}[AC-K29-irregular: coefficient-field vanishing at split primes]
\label{AC-K29-irregular}
Let \(p\geq5\), \(p\neq269\), split in \(K\).  Let \(f\)
and \(V=M_f(e)\), \(e\geq1\), satisfy the actual local,
non-CM and characteristic-zero derived-image hypotheses of
AC-K29-ramified.  No assumption
\(\chi(\mathfrak p)\neq\chi(\overline{\mathfrak p})\) is needed.
Then
\[
                 H_f^1(K,V\chi)=H_f^1(K,V\chi^2)=0.    \tag{CIR.15}
\]
This is a coefficient-field assertion, not an integral
regulator or all-prime size bound.
\end{theorem}

\begin{proof}
Apply the proved quadratic good twist, if necessary, to obtain
the good ordinary form \(f_0\), and restore the original
representation by the finite cyclotomic specialization as in
AC-K29-ramified.  Apply the universal ordinary Beilinson--Flach
class to the integral quotients (CIR.10), for the two CM
branches, before putting \(X=0\).  Their actual ray maps give
the same degree-two norm relation over \(K\) as AC-K29-ray;
the proof of that relation is the universal Hecke calculation,
which precedes this coefficient specialization.  The only
lattice denominators come from the fixed first-factor quotient
and the fixed primitive tame data.  They are independent of
the ray modulus.  The cyclotomic universal-norm argument gives
the Bloch--Kato conditions away from \(p\).

These universal maps do not require the full ordinary module
to be free over its Hecke algebra.  Precisely, KLZ,
Theorem 7.2.3(i),(ii),(v), gives weight-algebra projectivity of
\(F^\pm\), the dualizing description of \(F^-\), and the
Hecke-free rank-one description of \(F^+\); the extra
\(p\)-distinguished freeness theorem is 7.2.8.  The opening of
\S8.1 and Definition 8.1.1 use arbitrary ordinary Hida
families.  The complete proof of Proposition 8.1.7 uses
finite-level unramified \(\mathbb Z_p\)-free \(F^{--}\)
modules, compact inverse limits, and Lemma 8.1.5, so its
vanishing applies before our quotient.  Theorem 8.2.3 is
functorial for every unramified profinite module.  After the
ordinary twist used in the proof of Theorem 8.2.8, apply this
map to the actual rank-one quotient
\(F^-M(\mathbf f)^*\otimes F^+M(\mathbf g)^*\).
Its rank-one Perrin--Riou interpolation is the formula in
that proof.  This does not identify the entire possibly
nonfree weight-one specialization of \(M(\mathbf g)^*\)
with the usual two-dimensional classical cohomology module.

At the base modulus \(\pi\) identifies the ordinary line
with the first induced coordinate, by (CIR.8)--(CIR.10).
Consequently the canonical ordinary functional from Ohta and
KLZ, Proposition 10.1.1(1), specializes nontrivially.  There is
no remaining factor of \(X\) and no compensating functional
pole to postulate.  The exact normalization deserves emphasis.
In the complete proof of that proposition, printed pp.~50--51,
\(\omega_{\mathbf g}\) is the canonical isomorphism
\[
 D\bigl(F^+M(\mathbf g)^*(-1-\mathbf k-\varepsilon_{\mathbf g})\bigr)
           \xrightarrow{\ \sim\ }\Lambda_{\mathbf g}^{\rm cusp}.
                                                               \tag{CIR.16}
\]
It is obtained by taking \(\operatorname{Hom}_{\Lambda_D}(-,\Lambda_D)\)
of Ohta's cusp-form isomorphism and using Hida's pairing
\((T,F)\mapsto a_1(TF)\).  Thus the Dieudonne-module ordinary
generator corresponding to \(1\) pairs with the normalized
theta eigenfamily for a branch \(\phi\) as
\(\phi(1)=a_1(F_\phi)=1\).
Base change of (CIR.16) along \(\phi\) is still an isomorphism
to the branch coefficient ring.  Transporting it across the
isomorphism \(\pi|_{F^+}\) remains an isomorphism and hence is
nonzero at weight one.  This uses the quotient-dual evaluation;
it does not project a theta eigenvector in
\(\omega_T\otimes_T A\), where a congruence factor could occur.
The congruence-ideal map in KLZ is \(\eta_{\mathbf f}\),
part (2) of Proposition 10.1.1, on the other, already fixed
factor.  Its arithmetic-point normalization was checked in
AC-K29-ray and is unchanged here.

Apply the actual ordinary regulator
identity KLZ, Theorem 10.2.2, with its weight-one scalar
interpolation from \S7.7 immediately after Theorem 7.7.2.
The primitive fixed tame repair is AC-K29-primitive-reciprocity.
The resulting scalar is a nonzero multiple of
\(L^{\rm prim}(f,\theta_\chi,e)\), which is nonzero by the
primitive complex nonvanishing already used there.  All
unramified or quadratic ramified local logarithm and Euler
factors are the nonzero factors checked in AC-K29-lower and
AC-K29-ramified.  Equality of the two CM unit roots does not
create any new factor in those formulas.

The two branch filtrations specialize to the different global
\(\chi\) and \(\chi^c\) lines.  Even when both restrict to the
same unramified character at \(p\), they are different
subspaces of the two-dimensional Artin representation, as
the fixed global separator proves.  KLZ, Proposition 8.1.7,
annihilates the corresponding \(F^{--}\)-component for each
actual universal class.  The nonzero dual exponential therefore
lies on the associated one-dimensional singular line; the two
lines are independent and span the two-dimensional singular
space.  Local Poitou--Tate duality forces the Bloch--Kato group
to be strict at both primes of \(K\) over \(p\).
One ray Euler system has a nonzero base class.  The
characteristic-zero rank-two Euler-system bound, with the
retained derived-image hypotheses and the fixed lattice
denominator just proved, kills that dual strict group.
This proves (CIR.15), by the same comparison as in
AC-K29-lower.  No inference about an integral index is made.
\end{proof}
